%%%%%%%%%%%%%%%%%%%%%%%%%%%%%%%%%%%%%%%%%%%%%%%%%%%%%%%%%%%%%%%%%%%%%%%%%%%%%%%%%%%%%%%%%%%%%%%%
%% Discussion, limitations and future work. Written 28 November 2026.
%%
%% The sample-size figures of sec:future come from runs/task_level/per_slot.csv:
%% n_required = n_units * (mde80 / 0.03)^2, the usual inverse-square scaling of a standard error.
%% Median over the twenty-seven clinical slots is 480, range 256 to 1,187.
%%%%%%%%%%%%%%%%%%%%%%%%%%%%%%%%%%%%%%%%%%%%%%%%%%%%%%%%%%%%%%%%%%%%%%%%%%%%%%%%%%%%%%%%%%%%%%%%

\section{Discussion}
\label{sec:discussion}

\subsection{Does the encoding transfer?}
\label{sec:discussion-transfer}

Narrowly, and to two tasks. The published single-corpus result does reappear outside the corpus it
came from: on the clean ParkinsonHW sub-cohort and on Zhu the minimap arm beats the bare trace
under the same descriptor by $+0.086$ and $+0.109$, and the mechanism controls put that gain where
the encoding puts it --- destroying \emph{where} the painted value sits removes it, while painting a
channel with no information about the writer does not reproduce it. Both of those corpora prescribe
the figure, and in both the gain is carried by a single slot: the static spiral and the drawn
spiral, $+0.152$ and $+0.119$, while their companion slots give $+0.007$ and $+0.060$. A reader who
wants the one-sentence answer should have this one: the effect is real, it is localised to one task
in each of two corpora, and over the thirty tasks of this study it pools to zero.

Three qualifications belong in the same breath, because each of them was measured rather than
suspected. The margins reported for this recipe are in good part a change of descriptor: the
identical bare trace read by ImageNet features instead of by LPQ is worth $+0.154$ and $+0.265$ on
those two corpora, and under the backbone no corpus here gains more than $+0.014$ from the encoding.
Neither of the two superior slots survives Holm correction over the thirty-slot family, while the
one contrast that does survive is $+0.020$ --- inside the margin we declared as noise. And the
corpus that carries the largest raw margin of all, the full ParkinsonHW release at $+0.334$, is the
one whose capture metadata alone reaches $0.97$.

\subsection{The rule we registered does not hold at the task level}
\label{sec:discussion-rule}

The ordering that motivated this study --- the encoding helps where the task fixes the figure and
not where the produced shape carries the label --- holds across corpora and fails inside them. With
corpus fixed effects the task property is worth $-0.0090$ with $p = 0.391$ over an exact enumeration
whose attainable minimum is $0.00012$. This is the second time this study has asked that question
with the instrument that can answer it and received a no; the first was the within-corpus contrast on
DiaGraMo, declared with its direction and failed with the sign reversed.

We take the consequence rather than the escape. The between-corpus ordering is real in the sense
that it describes the six corpora, but what it tracks is not the task property alone: across corpora
the instrument, the recruitment, the population and the task all move together, and the component
that survives being isolated is not the task. The honest reading of the corpus-level moderator ---
coefficient $+0.098$, nine tenths of the heterogeneity absorbed, $p$ at its attainable minimum --- is
that it is a good description of six points, not a tested mechanism.

What replaces it is weaker and, we think, more useful. The two sides of the property differ in
\emph{dispersion} rather than in mean: the figure-imposed slots are heterogeneous around zero
($\hat\tau = 0.021$, $I^2 = 65\,\%$), the shape-informative slots homogeneous and nearly flat
($\hat\tau = 0.006$, $I^2 = 17\,\%$). Imposing the figure appears to be a necessary condition for a
large margin and nowhere near a sufficient one: ten of DiaGraMo's imposed slots produce nothing, and
three of them are reliably negative. That is a hypothesis for a study designed around it, and it is
stated here as one.

\subsection{Task-to-task variation exceeds corpus-to-corpus variation}
\label{sec:discussion-unit}

The quantitative result we did not expect is the variance decomposition. Between slots of one corpus
$\hat\tau = 0.0098$; between corpora $\hat\tau = 0.0067$. Both are below the median sampling
standard deviation of a single slot. A field that asks ``does representation $X$ work on dataset
$Y$'' is asking about the smaller of the two components, with an instrument whose noise exceeds
both.

This has a direct methodological consequence, and it is the one we would most like to be carried
forward: the task, not the corpus, is the natural unit for claims of this kind. It is also almost
free to adopt. The thirty estimates here were produced without rendering an image, fitting a model
or drawing a split; they are a re-slice of prediction files that any pipeline of this shape already
writes. The per-task breakdown is what shows that two corpus-level successes are two task-level
successes, which is not the same claim and not the one a reader of the corpus-level table would
infer.

\subsection{Equivalence was not available to be claimed}
\label{sec:discussion-equivalence}

We set out to test whether the encoded image is \emph{as good as} the bare trace where the shape
already carries the label, and the answer is that these corpora cannot say. Four of fifteen
shape-informative slots are equivalent at $\delta = 0.03$; eleven are inconclusive, and twenty-six
of the thirty slots overall have a minimum detectable effect larger than the margin itself. Three of
the four that do resolve belong to a non-clinical corpus with one writer.

We report this as a failure of H16 and not as a null finding, because the distinction is the point.
``No difference was found'' and ``the difference was shown to be smaller than $0.03$'' are different
claims, and only the second is a result. The earlier framing of this work made the first kind of
statement repeatedly, in the form of margins declared to be inside a $\pm 0.03$ band; every one of
those was a statement about our precision rather than about the representation.

\subsection{What a reader should take to the next cohort}
\label{sec:discussion-practice}

Five things, none of them expensive, and all of them things this study got wrong before it got them
right. Report the corpus's acquisition floor, so a score has a reference point; two of the corpora
here are partly or wholly separable from capture metadata, and one of them is the corpus the
original result came from. Give each descriptor its own bare-trace baseline, so a margin is about
the channel and not about the features. Report what effect the unit could have detected, so that a
null can be read at all. Measure the split noise instead of inheriting a band. And break the
endpoint down by task, because a corpus-level margin can be one task's margin diluted by seven
others --- or, as on DiaGraMo, a mixture of positive and reliably negative slots that averages to
nothing.

\subsection{Limitations}
\label{sec:limitations}

Thirty slots are not thirty independent units: fourteen come from one corpus and eight from
another, so the within-corpus component of the variance decomposition is estimated mainly inside
DiaGraMo, and the restricted permutation draws its variation from only two corpora. The
Drot\'ar--Dobe\v{s} recordings carry no task boundaries, so the second dysgraphia cohort contributes
one slot and cannot inform the within-corpus contrast at all. The acquisition floor is carried as a
corpus attribute rather than re-estimated per slot, which is a measurement this layer does not have.
Character Trajectories has a single writer, so its three slots support no participant-level claim
and enter as recording-level units. PaHaW's task subset was chosen in earlier work, so its eight
slots are not blind; the primary and the PaHaW-free analyses agree, which is why the verdicts stand.
The moderator's direction was informed by our own corpus-level results, as the registration states.
Pooling slots weights a corpus by how many tasks it administered, which makes the slot-level and
corpus-level pooled estimates two estimands rather than two estimates. And the selected
representation is chosen inside each corpus's training folds, not inside a slot, so a slot's margin
is the corpus's chosen arm evaluated there --- deliberately, since selecting per slot would widen
the family, but it means a slot-specific encoding might do better than any number here.

Finally, the screening of \cref{sec:screening} is a limitation as much as a method. Nineteen of
twenty-seven candidate corpora could not be obtained by download, and the eight that could are not a
random sample of handwriting research; they are the corpora whose authors chose open deposition.

\subsection{Future work}
\label{sec:future}

\paragraph{Corpora sized for the question.} The most actionable number in this paper is a design
number. Scaling each slot's standard error by the usual inverse square root, resolving a margin of
$0.03$ at $80\,\%$ power would take a median of $480$ participants per task across the twenty-seven
clinical slots here, with a range from $256$ to $1{,}187$; seventeen of those twenty-seven would
need more than twice the participants they have. The largest clinical corpus in this study has
$273$. Until collections of that size exist, work in this area can report effects and should not
report equivalence, and a study that wants to compare two representations on one clinical task needs
to say in advance what effect it could see.

\paragraph{A corpus built for the contrast.} The scarce resource in this design is a corpus that
elicits both kinds of task from the same participants on the same device in the same session. There
are two here, and one of them contributes a single figure-imposed slot. A prospective collection
that administers a balanced battery --- the same participants, half imposed figures and half free
content --- would test the task property with the confound held fixed by construction rather than by
a fixed effect, and it is the one experiment that could overturn the verdict of
\cref{sec:stats-moderator}.

\paragraph{A floor per task.} The acquisition floor is currently a corpus attribute. Computing it
per slot would show whether the capture artefacts that make two of these corpora partly separable
are spread across their battery or concentrated in particular tasks, which would change how a
contaminated corpus is salvaged: a corpus with a clean subset of tasks is more useful than a corpus
discarded whole.

\paragraph{Task boundaries in the unsegmented corpus.} Locating the elicited tasks inside the
Drot\'ar--Dobe\v{s} recordings would turn one slot into five and would be the first independent test
of the task-level contrast outside DiaGraMo. It must be done as its own study, with the
segmentation fixed before any score is computed, because a segmentation tuned against a metric would
compromise every number that depends on it.

\paragraph{Selection inside the task.} The arms here are the corpus's chosen encoding evaluated per
slot. Whether choosing the encoding \emph{per task}, inside training folds, recovers more than the
corpus-level choice is a well-posed question with a declared family size, and it is the one
extension of the method itself that these results suggest is worth the multiplicity it costs.
